% ===========================================================================
%  ALGTG — Active Learning through the Graph Transduction Game
%  Generic arXiv preprint (self-contained: compiles with a stock TeX Live /
%  MacTeX install, no conference style file required).
%
%  BUILD:  pdflatex main -> bibtex main -> pdflatex main -> pdflatex main
%
%  STATUS: Related Work, Background and Method are drafted. Everything marked
%  %% TODO needs your real numbers / figures after the baselines are re-run.
% ===========================================================================

\documentclass[11pt]{article}

\usepackage[utf8]{inputenc}
\usepackage[T1]{fontenc}
\usepackage[margin=1in]{geometry}
% lmodern rather than `times`: it ships complete T1 metrics, so the source also compiles on a
% minimal TeX installation (BasicTeX), where Times pulls in Courier metrics that are absent.
\usepackage{lmodern}
\usepackage{amsmath}
\usepackage{amssymb}
\usepackage{amsthm}
\newtheorem{proposition}{Proposition}
\usepackage{graphicx}
\usepackage{booktabs}
\usepackage{caption}
\usepackage[round]{natbib}          % author-year citations via \citep / \citet
\usepackage{url}
\usepackage[colorlinks=true,linkcolor=blue,citecolor=blue,urlcolor=blue]{hyperref}

\title{\textbf{Active Learning through the Graph Transduction Game}}
\author{
  Riccardo Zuliani\thanks{Correspondence: \texttt{riccardo.zuliani99@gmail.com}} \\
  %% TODO: affiliation and co-authors
  \textit{Affiliation}
}
\date{}

\begin{document}
\maketitle

% ---------------------------------------------------------------------------
\begin{abstract}
%% TODO — finalize once results are in. Draft:
Deep networks are data-hungry, yet high-quality annotation is expensive, which
renews the case for active learning (AL): selecting the few samples whose labels
help the most. We revisit AL through the lens of the \emph{Graph Transduction Game}
(GTG), an evolutionary, non-cooperative game whose replicator dynamics propagate
labels from a small labelled set to unlabelled points and settle at a consistent
(Nash) labelling. Our key observation is that the \emph{trajectory} of a sample's
label entropy under these dynamics is a strong signal of informativeness: points
that remain uncertain or oscillate the longest are the most valuable to annotate.
We introduce \textbf{ALGTG}, an acquisition function built on this signal. We further
show that avoiding redundant queries need not be hard-coded into the payoff: once
labelled players are clamped, large parts of the payoff matrix provably do not
influence the game, and the suppression of samples already well explained by the
labelled set \emph{emerges} from the dynamics themselves. We make this precise and
add a single-parameter term that amplifies the effect where it is actually
effective. Finally, we present an end-to-end differentiable variant that learns the
signal jointly with the backbone. On CIFAR-10/100, SVHN, Fashion-MNIST and
Tiny-ImageNet, ALGTG %% TODO: headline result vs. strongest baseline.
\end{abstract}

% ---------------------------------------------------------------------------
\section{Introduction}

Modern deep networks owe much of their success to large annotated corpora, but
annotation is the one ingredient that does not become cheaper with scale. In
specialised domains---medical imaging, remote sensing, speech---labels require
expert time, and the marginal value of an additional label varies enormously from
sample to sample. Active learning (AL) addresses exactly this asymmetry: given a
fixed annotation budget, choose the samples whose labels will improve the model the
most \citep{settles2009active}.

Most deep AL methods answer that question with a \emph{static} score. Uncertainty
methods read a scalar off the current model---predictive entropy, margin, or the
mutual information between prediction and parameters
\citep{houlsby2011bayesian,gal2017deep}. Diversity methods ignore the model's confidence and instead cover the
representation geometrically \citep{sener2018active}. Hybrids combine the two cues in
a single embedding or objective \citep{ash2020deep,parvaneh2022active}, and a further
line trains an auxiliary head to predict the signal
\citep{yoo2019learning,kim2021task,kye2023tidal}. What these approaches share is that
the score is a snapshot: it is computed from a frozen model at one instant, and it
says nothing about how a hypothetical label assignment would \emph{behave} once it
had to be reconciled with the labels of every other sample.

Transductive learning offers precisely that missing view. In the Graph Transduction
Game (GTG) of \citet{erdem2012graph}, labelling is cast as a non-cooperative game:
samples are players, labels are strategies, and a payoff rewards agreeing with
similar players. The replicator dynamics---a continuous generalisation of relaxation
labelling \citep{pelillo1997dynamics}---propagate the labelled seeds through the data
and settle at a Nash equilibrium, a labelling consistent with both the supervision and
the manifold structure. Existing work uses the \emph{outcome} of this process: the
equilibrium labelling. Our starting point is that the \emph{process itself} is
informative, and that for active learning it is in fact the more useful object.

The reason is intuitive. A sample lying deep inside a region the labelled set already
explains receives a strong, coherent pull toward one label and commits after a couple
of iterations. A sample on a decision boundary receives conflicting pulls and stays
ambivalent; a sample in a region no label reaches receives almost no pull at all and
remains undecided. Tracking each sample's label entropy \emph{across} the iterations
therefore yields a single quantity in which uncertainty and redundancy are already
entangled the way an acquisition function needs them to be---without hand-combining
two heuristics through a trade-off hyperparameter. This is the signal ALGTG queries on.

Making this precise required understanding which parts of the game actually matter.
Because labelled players are clamped to their known labels, we show
(Proposition~\ref{prop:inert}) that entire blocks of the payoff matrix cannot
influence the outcome: the natural-looking idea of making already-annotated samples
``repel'' each other, so as to discourage redundant queries, provably does nothing. The
redundancy avoidance that such a term is meant to produce is instead an
\emph{emergent} property of the dynamics, and we give a one-parameter term that
modulates it where it is genuinely effective. Finally, since the whole procedure is
differentiable, we derive an end-to-end variant in which the acquisition signal is
learned jointly with the backbone, letting representation and selection co-adapt.
%% TODO: closing sentence with the empirical headline once results are in.

\noindent\textbf{Contributions.}
\begin{itemize}
  \item A novel AL acquisition function based on the \emph{entropy trajectory} of
        samples under Graph Transduction Game dynamics.
  \item An analysis of \emph{which blocks of the payoff matrix actually drive the
        game} once labelled players are clamped (Proposition~\ref{prop:inert}). It
        follows that redundancy avoidance is an emergent property of GTG rather than
        something that must be engineered into the payoff; we also derive a
        one-parameter term that amplifies it where it is effective.
  \item An end-to-end differentiable variant that learns the acquisition signal
        jointly with the task network, so representation and selection co-adapt.
  \item %% TODO: empirical claim — SOTA / competitive on N benchmarks with M trials.
\end{itemize}

% ---------------------------------------------------------------------------
\section{Related Work}

\paragraph{Uncertainty-based AL.} Classical strategies query the points where the
current model is least confident, using predictive entropy, margin or least
confidence \citep{settles2009active}. Bayesian variants estimate epistemic
uncertainty via the mutual information between predictions and model parameters
(BALD) \citep{houlsby2011bayesian}, made practical for deep nets through MC-dropout
\citep{gal2017deep} and extended to the batch setting by BatchBALD
\citep{kirsch2019batchbald}. Pure uncertainty ignores redundancy and tends to
select correlated batches.

\paragraph{Diversity and hybrid AL.} Representation-driven methods instead cover
the feature space: Core-Set casts batch selection as a $k$-center problem over
embeddings \citep{sener2018active}. Hybrid methods combine both cues: BADGE runs
$k$-means\texttt{++} on \emph{gradient} embeddings, whose magnitude encodes
uncertainty and direction encodes diversity \citep{ash2020deep}; AlphaMix probes
for decision-boundary instability by interpolating unlabelled features toward
labelled class anchors \citep{parvaneh2022active}; CDAL measures contextual
diversity through KL divergence of predictive distributions \citep{agarwal2020contextual}.

\paragraph{Learned acquisition.} A parallel line learns the acquisition signal.
Learning Loss attaches a module that predicts each sample's loss and queries the
highest-loss points \citep{yoo2019learning}. VAAL trains a VAE--discriminator to
tell labelled from unlabelled in latent space \citep{sinha2019variational}, and
TA-VAAL makes it task-aware by feeding in a learned loss/rank
\citep{kim2021task}. TiDAL exploits training dynamics as a supervisory signal for
the acquisition head \citep{kye2023tidal}.

\paragraph{Graph-based and transductive AL.} Closest to our work, Sequential GCN
builds a graph over labelled and unlabelled samples and learns node
representations to rank informativeness \citep{caramalau2021sequential}. Our method
is also graph-based but differs fundamentally: rather than learning graph
embeddings, we run the \emph{Graph Transduction Game}
\citep{erdem2012graph}---an evolutionary game whose replicator dynamics generalize
classical relaxation labelling \citep{pelillo1997dynamics}---and derive the
acquisition signal from the \emph{temporal evolution} of label entropy under those
dynamics. To our knowledge, GTG dynamics have not previously been used as an AL
acquisition function.

% ---------------------------------------------------------------------------
\section{Background: The Graph Transduction Game}

Given $n$ objects (players) and $c$ labels (pure strategies), GTG casts
transductive labelling as a non-cooperative game \citep{erdem2012graph}. Player
$i$ holds a mixed strategy $x_i \in \Delta^{c-1}$ (a distribution over labels),
and players receive payoff for agreeing with \emph{similar} players, encoded by a
non-negative affinity matrix $W \in \mathbb{R}^{n\times n}$. Writing
$x=(x_1,\dots,x_n)$, the population evolves by the multi-population replicator
dynamics
\begin{equation}
  x_{i}^{(t+1)} \;=\; x_{i}^{(t)} \odot
  \frac{\big(W x^{(t)}\big)_{i}}{\,x_{i}^{(t)\top}\big(W x^{(t)}\big)_{i}\,},
  \label{eq:replicator}
\end{equation}
where $\odot$ is the elementwise product and the ratio is taken per label. Labelled
players are clamped to their one-hot strategies throughout. The dynamics increase
average payoff monotonically and converge to a Nash equilibrium---a labelling that
is consistent with the affinity structure and the labelled seeds.

% ---------------------------------------------------------------------------
\section{Method: ALGTG}

Let $L$ and $U$ be the labelled and unlabelled sets and $\phi(\cdot)$ the backbone
feature map. For a pool of $n=|L|+|U|$ samples we build a payoff matrix, run the
game, and score unlabelled samples by their entropy trajectory.

\paragraph{Payoff construction.} We form $W$ from a similarity kernel over features,
$W_{ij}=k\big(\phi(x_i),\phi(x_j)\big)$, with $k$ one of cosine similarity,
Pearson correlation, or a self-tuning RBF kernel whose per-point bandwidth is the
distance to the $7$th nearest neighbour \citep{zelnikmanor2004self}. Because the
replicator update~\eqref{eq:replicator} requires a non-negative payoff to keep
$x_i$ on the simplex, we clamp $W$ to $[0,1]_{\ge 0}$ and use a zero diagonal, so no
player is rewarded for agreeing with itself. An optional threshold (a fixed value or
the matrix mean) sparsifies weak edges.

\paragraph{Initialization and dynamics.} We set $x_i^{(0)}=e_{y_i}$ for $i\in L$
and $x_i^{(0)}=\tfrac{1}{c}\mathbf{1}$ for $i\in U$, then iterate
\eqref{eq:replicator} re-clamping the labelled rows each step, until
$\lVert x^{(t+1)}-x^{(t)}\rVert < \varepsilon$ or a maximum of $T$ iterations. If a
player receives no payoff from anyone---an isolated node, which sparsification can
create---its update is $0/0$; we adopt the natural convention that such a player has
no incentive to deviate and keep $x_i^{(t+1)}=x_i^{(t)}$. This keeps every player on
the simplex; normalizing by a floored denominator would instead collapse the row to
the zero vector, silently removing the point from the game and assigning it zero
entropy.

\paragraph{Which blocks of the payoff drive the game?} Partitioning $W$ into blocks
$W_{LL},W_{LU},W_{UL},W_{UU}$ induced by $L$ and $U$ yields a simple but consequential
observation.

\begin{proposition}\label{prop:inert}
If the labelled players are clamped at every iteration, then the trajectory
$\{x_U^{(t)}\}_{t\ge 0}$ of the unlabelled players---and hence any acquisition score
derived from it---depends on $W$ only through $W_{UU}$ and $W_{UL}$. The blocks
$W_{LL}$ and $W_{LU}$ are inert.
\end{proposition}

\begin{proof}[Proof sketch]
The update of player $i$ in \eqref{eq:replicator} depends on $W$ only through row
$i$, namely $(Wx)_i=\sum_j W_{ij}x_j$. For $i\in U$ this row meets only the columns
of $L$ and $U$ within $W_{UL}$ and $W_{UU}$. Rows indexed by $i\in L$---the only rows
in which $W_{LL}$ and $W_{LU}$ appear---are overwritten by the clamp before they are
used. \end{proof}

Proposition~\ref{prop:inert} has a practical corollary: any attempt to steer the
acquisition by editing $W_{LL}$ (for instance, turning labelled--labelled affinities
into distances so that annotated points ``repel'' each other, an intuitive design we
initially adopted) has \emph{no effect whatsoever} on the selected batch. We verify
this numerically: such edits leave the acquisition scores identical to machine
precision.

\paragraph{Redundancy avoidance is emergent.} Rather than a defect, this delimits
where the mechanism actually lives. An unlabelled point lying close to the labelled
set receives, through $W_{UL}$, a strong and coherent pull toward a single label; it
commits almost immediately, so its entropy trajectory is short and its score low.
A point in a region the labelled set does not explain receives weak or conflicting
pull and remains uncertain for many iterations, earning a high score. Standard GTG
dynamics therefore already suppress redundant queries, with no bespoke payoff term.
To \emph{modulate} the strength of this effect we introduce a single non-negative
parameter $\lambda$ that scales the labelled--unlabelled coupling,
\begin{equation}
  W_{UL} \leftarrow (1+\lambda)\,W_{UL},
  \qquad W_{LU} \leftarrow (1+\lambda)\,W_{LU},
  \label{eq:redundancy}
\end{equation}
with $\lambda=0$ recovering standard GTG. Because clamped labelled players enter the
potential of the restricted game on $U$ only through a linear term, \eqref{eq:redundancy}
preserves non-negativity and the symmetry of $W_{UU}$, so the Baum--Eagon convergence
guarantee is unaffected. Empirically, increasing $\lambda$ monotonically widens the
score gap between redundant and informative samples. %% TODO: cite the ablation figure.

\paragraph{Acquisition by entropy trajectory.} For each unlabelled sample we record
the Shannon entropy of its strategy at every iteration,
$H_i^{(t)} = -\sum_{k} x_{ik}^{(t)} \log x_{ik}^{(t)}$, and aggregate the trajectory
either by its mean or by its area (trapezoidal integral):
\begin{equation}
  s_i \;=\; \tfrac{1}{T}\!\sum_{t=1}^{T} H_i^{(t)}
  \quad\text{or}\quad
  s_i \;=\; \textstyle\int_{1}^{T} H_i^{(t)}\,dt .
\end{equation}
We query the top-$k$ unlabelled samples by $s_i$. Intuitively, a point whose label
assignment settles immediately is easily explained by its neighbours and carries
little new information, whereas a point that stays uncertain or oscillates under
propagation lies where the current model and manifold structure disagree---exactly
where a label is most useful.

% Typeset without the algorithm/algorithmic packages so the source compiles on a minimal
% TeX installation (e.g. BasicTeX) as well as on a full TeX Live.
\begin{figure}[t]
\hrule\vspace{3pt}
\noindent\textbf{Algorithm 1}\quad ALGTG acquisition (offline)
\vspace{3pt}\hrule\vspace{5pt}
\begin{tabbing}
\hspace{1.6em}\=\hspace{1.4em}\=\kill
\textbf{Require:} features $\phi(x)$ for $L\cup U$; budget $k$; iterations $T$; tolerance $\varepsilon$ \\[2pt]
1: \> build the payoff $W$ (kernel, clamp to $[0,1]_{\ge0}$, zero diagonal, threshold) \\
2: \> initialise $x^{(0)}$: one-hot for $i\in L$, uniform for $i\in U$ \\
3: \> \textbf{for} $t=0,\dots,T-1$ \textbf{do} \\
4: \> \> update $x^{(t+1)}$ by Eq.~\eqref{eq:replicator}; re-clamp the labelled rows \\
5: \> \> record $H_i^{(t+1)}$ for $i\in U$ \\
6: \> \> \textbf{if} $\lVert x^{(t+1)}-x^{(t)}\rVert<\varepsilon$ \textbf{then break} \\
7: \> \textbf{end for} \\
8: \> $s_i \leftarrow \mathrm{aggregate}\big(H_i^{(1:T)}\big)$ for $i\in U$ \\
9: \> \textbf{return} the top-$k$ unlabelled samples ranked by $s_i$
\end{tabbing}
\vspace{-4pt}\hrule
\label{alg:algtg}
\end{figure}

\paragraph{End-to-end differentiable variant.} The offline scheme runs GTG on
frozen embeddings after training. We also propose an online variant in which a
lightweight module projects intermediate backbone features to a low-dimensional
space, runs the \emph{differentiable} dynamics~\eqref{eq:replicator} on the induced
affinity, and reads the entropy trajectory through an MLP or LSTM head to predict
$s_i$. The head is trained jointly with the task loss by regressing the GTG entropy
computed on the full embeddings (a teacher signal), so gradients from the
acquisition objective also shape the representation. This mirrors the auxiliary
training of Learning Loss \citep{yoo2019learning} but with a game-theoretic,
transduction-derived target. %% TODO: state which variant is the final "ALGTG"
%% and move the others to the ablation.

% ---------------------------------------------------------------------------
\section{Experiments}

\paragraph{Datasets.} We evaluate on five image-classification benchmarks spanning a range
of label cardinalities and image statistics: CIFAR-10 and CIFAR-100 ($50$k train
/ $10$k test, $32\times32$), SVHN ($73$k / $26$k, $32\times32$), Fashion-MNIST
($60$k / $10$k, $28\times28$, greyscale) and Tiny-ImageNet ($100$k / $10$k, $64\times64$,
$200$ classes).

\paragraph{Protocol.} Every method starts from the same randomly drawn set of $1000$
labelled samples and performs $10$ acquisition cycles of $1000$ labels each, ending at
$11$k labels. At each cycle the candidate pool is a fresh random subset of $10{,}000$
unlabelled samples, which bounds the cost of the pool-based methods (ours included) and is
applied identically to every baseline. After each acquisition the network is
\emph{re-initialised from the same checkpoint and retrained from scratch} for $200$ epochs,
so that the reported accuracy reflects the acquired labels and not the optimisation history.
Each configuration is repeated over $5$ trials with different seeds, which also re-draw the
initial labelled set; we report the mean and a $95\%$ confidence interval.

\paragraph{Implementation.} The backbone is ResNet-18 \citep{he2016deep} (VGG-16-bn for
Caltech-256), trained with SGD (lr $0.1$, momentum $0.9$, weight decay $5\times10^{-4}$,
step decay at epoch $160$) and batch size $128$. The learnable acquisition modules use Adam
(lr $10^{-3}$). Unless stated otherwise ALGTG uses the correlation kernel, mean-thresholded
sparsification, mean aggregation of the entropy trajectory, $T=30$ iterations, tolerance
$\varepsilon=10^{-4}$ and $\lambda=0$.

\paragraph{Metrics.} Alongside accuracy at each budget we report the area under the
budget--accuracy curve (AUBC), normalised by the budget range, which summarises a method's
behaviour over the whole annotation trajectory in a single number and avoids conclusions that
hinge on one cherry-picked budget.

\paragraph{Reproducibility.} All methods share the same backbone, optimiser, schedule, seeds
and candidate pools, and differ only in the acquisition function, so the comparison isolates
the selection strategy. Every baseline is re-implemented in the same framework rather than
copied from published tables. We flag one consequence of this explicitly: while preparing the
study we found defects in three of our own baseline implementations---MC-dropout that produced
no stochasticity (so BALD's mutual information was identically zero), per-class anchoring that
silently fell back to a global mean in AlphaMix, and BADGE clustering penultimate features
instead of gradient embeddings---each of which had made the corresponding baseline weaker than
it should be. All reported numbers were produced after these were fixed; no pre-fix result is
included. Code, configurations and the scripts that turn the raw logs into the tables and
figures of this paper are released with it.

\paragraph{Baselines.} Random, Entropy, Core-Set \citep{sener2018active},
BADGE \citep{ash2020deep}, BALD \citep{gal2017deep}, Learning Loss
\citep{yoo2019learning}, VAAL/TA-VAAL \citep{sinha2019variational,kim2021task},
CDAL \citep{agarwal2020contextual}, AlphaMix \citep{parvaneh2022active},
TiDAL \citep{kye2023tidal}, and both variants of the graph-based Sequential GCN
---UncertainGCN and CoreGCN---\citep{caramalau2021sequential}, which are the closest
competitors to our setting.
%% NOTE: BALD, AlphaMix and BADGE were re-run after correctness fixes; do not
%% report any pre-fix numbers.

% Main results table. Generated from the benchmark CSV by
%     python docs/paper/make_results.py results/<exp_path> cifar10
% which writes results_<dataset>.tex next to this file. Until that has been run the
% placeholder below is used, so the document always compiles.
\IfFileExists{results_cifar10.tex}{\input{results_cifar10}}{%
  \begin{table}[t]
  \centering
  \caption{Test accuracy (\%) vs.\ labelling budget. \emph{Placeholder --- run
  \texttt{docs/paper/make\_results.py} to populate.}}
  \label{tab:main-cifar10}
  \begin{tabular}{lccc}
  \toprule
  Method & 2k & 5k & 10k \\
  \midrule
  Random                & -- & -- & -- \\
  \dots                 & -- & -- & -- \\
  \textbf{ALGTG} (ours) & -- & -- & -- \\
  \bottomrule
  \end{tabular}
  \end{table}%
}

\paragraph{Ablations.} %% TODO — each isolates one design choice:
(i) affinity kernel (cosine / correlation / RBF);
(ii) entropy aggregation (mean vs.\ integral);
(iii) offline vs.\ online, and head architecture (MLP vs.\ LSTM);
(iv) the redundancy weight $\lambda$ of Eq.~\eqref{eq:redundancy}, swept from $0$
     (standard GTG) upward, quantifying how much explicit amplification adds over
     the emergent behaviour;
(v) \textbf{the central control}: the acquisition signal itself---the full entropy
     \emph{trajectory} versus only the entropy at convergence, versus the plain
     softmax entropy of the backbone. This isolates whether the \emph{dynamics},
     and not merely the final labelling or the model's own confidence, carry the
     information.

\paragraph{Analysis.} %% TODO — entropy-trajectory plots separating queried vs.\
non-queried samples; t-SNE of selections vs.\ competitors; cost/runtime discussion.

% ---------------------------------------------------------------------------
\section{Conclusion}

We presented ALGTG, an active-learning acquisition function that scores unlabelled
samples by the \emph{trajectory} of their label entropy under the replicator dynamics
of a graph transduction game, rather than by a static readout of a frozen model. Along
the way we characterised which parts of the game's payoff can influence the outcome at
all (Proposition~\ref{prop:inert}), which shows that the avoidance of redundant queries
is emergent rather than something to be engineered, and we gave a differentiable
formulation that lets the acquisition signal shape the representation it is computed
from. %% TODO: one sentence with the headline empirical result.

\paragraph{Limitations.} Three are worth stating plainly. First, \emph{cost}: the payoff
matrix is quadratic in the size of the candidate pool, so the offline variant becomes
expensive as the labelled set grows, and the online variant unrolls the dynamics inside
the autograd graph. Sub-sampling the pool, as we do, bounds this but makes the
transduction depend on the sample drawn. Second, \emph{locality of the online variant}:
it runs the game over a mini-batch rather than the full pool, so its equilibrium---and
hence its signal---depends on batch composition; it is best understood as an auxiliary
objective that shapes the backbone, in the spirit of \citet{yoo2019learning}, rather
than as a cheaper substitute for the offline computation, whose teacher signal is
available at query time anyway. Third, \emph{graph construction}: the acquisition
inherits whatever the affinity kernel and its sparsification encode, and aggressive
thresholding can disconnect points from the graph entirely, at which point the dynamics
leave them untouched. We keep such nodes on the simplex by convention, but a principled
treatment of disconnected candidates---which are often precisely the novel ones---is
left to future work.

% ---------------------------------------------------------------------------
\bibliographystyle{plainnat}
\bibliography{references}

\end{document}
